% =====================================================================
% 表1「主要符号及其含义」草稿（供 Paper/5.Symbol_Description/content.tex 替换用）
%
% 用法：把下面 \begin{table}...\end{table} 整块复制，
%       覆盖 content.tex 里现在那个只有 4 行占位符（【样本数量】等）的表。
%       本文件不被 main.tex 引用，不会影响编译，可随时删除。
%
% 说明：
%   1) 只收录「一~七」章出现的符号；论文三（问题三）与问题四的模型仍在修改中，
%      为避免变量表反复返工，待其定稿后再把新增符号追加到本表末尾即可。
%   2) 单位一列只写物理量单位；纯计数/无量纲量写「---」。
%   3) 表中每个符号在正文首次出现处也都有文字定义，本表是汇总索引。
% =====================================================================
\begin{table}[htbp]
  \centering
  \caption{主要符号及其含义}\label{tab:notation}
  \begin{tabularx}{0.95\linewidth}{cXc}
    \toprule
    符号 & 含义 & 单位 \\
    \midrule
    \multicolumn{3}{@{}l}{\textit{索引与集合}}\\
    \(t\) & 时段编号，\(t=1,2,\ldots,T\)，每日 \(T=144\) & --- \\
    \(\tau_t\) & 第 \(t\) 时段对应的时刻 & h \\
    \(\Delta t\) & 单位时段长度，\(\Delta t=1/6\) h & h \\
    \(d\) & 调度日编号（问题二逐日滚动） & --- \\
    \(m\) & 场景编号，\(m=1,2,\ldots,M\) & --- \\
    \(b_m\) & 场景 \(m\) 抽中的历史残差日期 & --- \\
    \midrule
    \multicolumn{3}{@{}l}{\textit{已知数据与参数}}\\
    \(\pi_t\) & 第 \(t\) 时段的外网购电价格 & 元/kWh \\
    \(L_t,\,G_t\) & 第 \(t\) 时段的实际负荷电量与光伏发电量 & kWh \\
    \(P_t^{\mathrm L},\,P_t^{\mathrm{PV}}\) & 第 \(t\) 时段的负荷功率与光伏预测功率 & kW \\
    \(P_t^{\mathrm{net}}\) & 第 \(t\) 时段的净负荷功率，\(P_t^{\mathrm L}-P_t^{\mathrm{PV}}\) & kW \\
    \(\hat L_t,\,\hat G_t\) & 第 \(t\) 时段的负荷与光伏点预测值 & kWh \\
    \(L_t^m,\,G_t^m\) & 场景 \(m\) 下第 \(t\) 时段的负荷与光伏发电量 & kWh \\
    \(D_t\) & 扣除储能放电后的净需求 & kWh \\
    \(q_{\max}\) & 单时段最大充放电量 & kWh \\
    \(S_{\min},\,S_{\max}\) & 储电量运行下界与上界 & kWh \\
    \(S_0\) & 初始（日初）储电量 & kWh \\
    \(\eta_c,\,\eta_r\) & 充电效率与放电效率 & --- \\
    \(\alpha,\,\beta\) & CVaR 的置信水平与风险权重 & --- \\
    \(\varepsilon\) & 第一阶段购电费用容差 & 元 \\
    \(\kappa_2\) & 软终端偏差惩罚系数 & --- \\
    \(\kappa_{2,\mathrm{base}}\) & \(\kappa_2\) 的缩放基准（问题一储能边际价值中位数） & --- \\
    \(M\) & 场景数 & --- \\
    \(\lambda\) & 软终端系数相对于基准值的倍数 & --- \\
    \midrule
    \multicolumn{3}{@{}l}{\textit{决策变量}}\\
    \(x_t\) & 第 \(t\) 时段的计划购电量 & kWh \\
    \(c_t,\,r_t\) & 第 \(t\) 时段的计划充电量与放电量 & kWh \\
    \(w_t\) & 第 \(t\) 时段的弃光量 & kWh \\
    \(s_t\) & 第 \(t\) 时段结束时的储电量 & kWh \\
    \(y_t^m,\,u_t^m\) & 场景 \(m\) 下的计划电提取量与未提取量 & kWh \\
    \(e_t^m\) & 场景 \(m\) 下的紧急购电量 & kWh \\
    \(g_t^m\) & 场景 \(m\) 下的实际光伏利用量 & kWh \\
    \(v_t^m\) & 场景 \(m\) 下的供给过剩弃电量 & kWh \\
    \(\xi^+,\,\xi^-\) & 日末储电量相对期初的正、负偏差 & kWh \\
    \(q_m\) & 场景 \(m\) 下超过风险阈值的费用部分 & 元 \\
    \midrule
    \multicolumn{3}{@{}l}{\textit{派生量与评价指标}}\\
    \(\zeta\) & 费用分布的风险阈值 & 元 \\
    \(C_m\) & 场景 \(m\) 下的购电费用 & 元 \\
    \(C_1^\ast\) & 第一阶段最低购电费用 & 元 \\
    \(C_{\mathrm{base}}\) & 无储能基准购电费用 & 元 \\
    \(C_{\mathrm{val}}\) & 参数验证期总费用 & 元 \\
    \(C_{\mathrm{emg}}\) & 紧急购电费用 & 元 \\
    \(v_t\) & 储能边际价值（储能状态约束对偶变量取负） & 元/kWh \\
    \(\mathrm{SOC}_t\) & 储电量占额定容量的比例 & \% \\
    \(\mathrm{WAPE}\) & 加权绝对百分比误差 & \% \\
    \(\mathrm{MAE}\) & 平均绝对误差 & kWh \\
    \bottomrule
  \end{tabularx}
\end{table}
